\section{Conclusion and Limitations}
\label{sec:conclusion}

%%% Conclusion 设计要点（正文已写，注释为背景） %%
% 风格：精炼，正文把 conclusion 与 limitations 合并成一段；回扣 intro 的三大卖点，不引入新内容。

%%% 1. 回扣任务与方法（走故事线，避免流水账式总结） %%%
% 动机：FM 范式自身两大临床缺陷——数据饥渴 + 实时性差 → 本文研究 = 修复二者
% 统一视角：手术视频双重冗余（空间背景 + 时间稳态帧）是数据缺陷背后的关键障碍
% 方法：对抗冗余两阶段——
%   修数据缺陷 SPR-Teacher：SSTS 压空间冗余 + KPTF 压时间冗余，冻结 generic FM + 少量标注；
%   修实时缺陷 SPR-Lite：CAP 跨架构 + STGD 时空先验稠密监督。
%   （不复述 generic vs surgical FM 对比——那已降级为引言一句 motivation）
%%% 2. 结果总结（方法 / 部署） %%%
% 方法：SPR-Teacher 在 4 公开 + 1 私有数据集 few-shot 大幅超越 prior SOTA。
% 部署：Lite 保留精度，参数大幅压缩、实时性显著提升。
%%% 3. 临床意义 + 未来工作 %%%
% 为低标注、实时可部署的智能手术室提供可行路径；
% 未来：更多术式与多模态、真实临床硬件验证、更强蒸馏目标。

%%% Limitation 设计要点（已并入同一段，正文已写，注释为背景） %%
% 总基调：坦诚、客观，笼统总结，不点名具体模型/数据集，不引用表格，落到"未来可做"上。

%%% 1. 只考虑公开权重的 FM，覆盖不全 %%%
% 后果：所用 FM 的全面性受限。
% 未来：期待社区开放更多权重。

%%% 2. surgical FM 预训练存在手术数据泄露，FM 对比未必公平 %%%
% 事实：surgical FM 不可避免地会在公开手术视频上预训练，部分可能与评测数据重叠、使其分数虚高，
%   故 generic vs surgical 的 FM 对比不完全公平。
% 未来：期待更严格的"未见术式"评估协议。

%%% 3. 实时性未在真实临床硬件上验证 %%%
% 未来：针对临床硬件做量化与工程优化。

Surgical phase recognition is a fundamental task for intelligent operating
rooms. In the foundation model era, its clinical deployment is hindered by
two deficiencies of the FM paradigm itself: a reliance on massive training
data and an inability to infer in real time on clinical hardware. We argue
that the dual redundancy of surgical videos---dominant background regions and
repetitive steady-state frames---is the key obstacle behind the data
deficiency, and propose to combat it in both the model and its
distillation. SPR-Teacher compresses spatial redundancy through salient spatial
token selection and temporal redundancy through keyframe-prior temporal
fusion, adapting a frozen generic foundation model with only a handful of
labeled videos. SPR-Lite then lifts the real-time deficiency by
compressing the model itself, where a
cross-attention projection bridges arbitrary student architectures and a
spatiotemporally-guided distillation transfers the teacher's where- and
when-to-look priors as dense supervision. Across four public and one private
dataset, SPR-Teacher surpasses prior methods under few-shot settings by large
margins, while SPR-Lite retains this performance with under 9M
parameters at over 25\,FPS on a single CPU core, offering a practical route
toward low-annotation, real-time deployment.
Nevertheless, several limitations remain. First, we only consider foundation
models with publicly released
weights, so the coverage may not be exhaustive. Second, surgical FMs
inevitably pre-train on public surgical videos, part of which may overlap with
our evaluation data and inflate their scores, making the FM comparison not
perfectly fair. Third, real-time efficiency is measured on a server CPU core
rather than on clinical hardware, where quantization and further engineering
optimization remain to be explored. We leave these for future work.

